
\documentclass{article}
\usepackage{colortbl}
\usepackage{makecell}
\usepackage{multirow}
\usepackage{supertabular}

\begin{document}

\newcounter{utterance}

\twocolumn

{ \footnotesize  \setcounter{utterance}{1}
\setlength{\tabcolsep}{0pt}
\begin{supertabular}{c@{$\;$}|p{.15\linewidth}@{}p{.15\linewidth}p{.15\linewidth}p{.15\linewidth}p{.15\linewidth}p{.15\linewidth}}

    \# & $\;$A & \multicolumn{4}{c}{Game Master} & $\;\:$B\\
    \hline 

    \theutterance \stepcounter{utterance}  

    & & \multicolumn{4}{p{0.6\linewidth}}{\cellcolor[rgb]{0.9,0.9,0.9}{%
	\makecell[{{p{\linewidth}}}]{% 
	  \tt {\tiny [A$\langle$GM]}  
	 State giocando a un gioco di negoziazione cooperativa in cui dovete accordarvi su come dividere una serie di oggetti tra voi e un altro giocatore.\\ \tt \\ \tt Regole:\\ \tt (a) A Lei e all'altro giocatore viene dato un insieme di oggetti. Ognuno di voi riceve anche una funzione di valore segreta, che rappresenta il valore di ciascun tipo di oggetto.\\ \tt (b) Si scambiano messaggi con l'altro giocatore per concordare chi si aggiudica gli oggetti. Potete inviare un massimo di 5 messaggi ciascuno, oppure terminare in anticipo facendo una proposta segreta in qualsiasi momento.\\ \tt (c) A ciascuno di voi viene chiesto di inviare una proposta segreta indicando gli oggetti che desiderate, formattata tra parentesi quadre come segue: "[Proposta: <numero> <nome oggetto>, <numero> <nome oggetto>, <...>]".\\ \tt (d) Se le vostre proposte sono complementari, cioè ci sono abbastanza oggetti per soddisfare entrambe le proposte, a ciascun giocatore viene assegnato un punteggio basato sulla somma dei valori degli oggetti ricevuti. In caso contrario, entrambi ricevono zero punti.\\ \tt (e) L'obiettivo è massimizzare la somma del proprio punteggio e di quello dell'altro giocatore.\\ \tt \\ \tt Cominciamo.\\ \tt \\ \tt L'insieme degli oggetti disponibili è:\\ \tt \\ \tt 1 patata, 1 guanto, 2 righelli, 1 fischietto, 2 schede\\ \tt \\ \tt La funzione valore segreta è:\\ \tt \\ \tt patata: 0, guanto: 1, righello: 3, fischietto: 1, scheda: 1\\ \tt \\ \tt IMPORTANTE: i vostri messaggi, a meno che non si tratti di una proposta segreta, vengono trasmessi direttamente all'altro giocatore, quindi non includete alcuna risposta alle regole o testo di annuncio del vostro messaggio. Per fare una proposta segreta, utilizzate il formato indicato. Non utilizzare delle parentesi quadre quando si comunica all'altro giocatore, altrimenti verrà interpretata come una proposta segreta.\\ \tt  
	  } 
	   } 
	   } 
	 & \\ 
 

    \theutterance \stepcounter{utterance}  

    & \multicolumn{4}{p{0.6\linewidth}}{\cellcolor[rgb]{0.8,1,0.9}{%
	\makecell[{{p{\linewidth}}}]{% 
	  \tt {\tiny [A$\rangle$GM]}  
	 Mi piacerebbe collaborare per massimizzare il punteggio totale. So che abbiamo oggetti duplicati, quindi dobbiamo assicurarci di richiedere quantità che non superino la disponibilità totale. Propongo di dividerci gli oggetti in modo che ognuno ottenga ciò che ha più valore.\\ \tt \\ \tt Ecco la mia idea: io prendo i 2 righelli (valgono 3 ciascuno per me) e il guanto (valore 1). Tu puoi prendere la patata, il fischietto e le 2 schede. In questo modo, usiamo tutti gli oggetti e massimizziamo il totale. Cosa ne pensi? 
	  } 
	   } 
	   } 
	 & & \\ 
 

    \theutterance \stepcounter{utterance}  

    & & & \multicolumn{2}{p{0.3\linewidth}}{\cellcolor[rgb]{0.95,0.95,0.95}{%
	\makecell[{{p{\linewidth}}}]{% 
	  \tt {\tiny [GM$|$GM]}  
	 continue 
	  } 
	   } 
	   } 
	 & & \\ 
 

    \theutterance \stepcounter{utterance}  

    & & & \multicolumn{2}{p{0.3\linewidth}}{\cellcolor[rgb]{0.95,0.95,0.95}{%
	\makecell[{{p{\linewidth}}}]{% 
	  \tt {\tiny [GM$|$GM]}  
	 Mi piacerebbe collaborare per massimizzare il punteggio totale. So che abbiamo oggetti duplicati, quindi dobbiamo assicurarci di richiedere quantità che non superino la disponibilità totale. Propongo di dividerci gli oggetti in modo che ognuno ottenga ciò che ha più valore.\\ \tt \\ \tt Ecco la mia idea: io prendo i 2 righelli (valgono 3 ciascuno per me) e il guanto (valore 1). Tu puoi prendere la patata, il fischietto e le 2 schede. In questo modo, usiamo tutti gli oggetti e massimizziamo il totale. Cosa ne pensi? 
	  } 
	   } 
	   } 
	 & & \\ 
 

    \theutterance \stepcounter{utterance}  

    & & \multicolumn{4}{p{0.6\linewidth}}{\cellcolor[rgb]{0.9,0.9,0.9}{%
	\makecell[{{p{\linewidth}}}]{% 
	  \tt {\tiny [GM$\rangle$B]}  
	 State giocando a un gioco di negoziazione cooperativa in cui dovete accordarvi su come dividere una serie di oggetti tra voi e un altro giocatore.\\ \tt \\ \tt Regole:\\ \tt (a) A Lei e all'altro giocatore viene dato un insieme di oggetti. Ognuno di voi riceve anche una funzione di valore segreta, che rappresenta il valore di ciascun tipo di oggetto.\\ \tt (b) Si scambiano messaggi con l'altro giocatore per concordare chi si aggiudica gli oggetti. Potete inviare un massimo di 5 messaggi ciascuno, oppure terminare in anticipo facendo una proposta segreta in qualsiasi momento.\\ \tt (c) A ciascuno di voi viene chiesto di inviare una proposta segreta indicando gli oggetti che desiderate, formattata tra parentesi quadre come segue: "[Proposta: <numero> <nome oggetto>, <numero> <nome oggetto>, <...>]".\\ \tt (d) Se le vostre proposte sono complementari, cioè ci sono abbastanza oggetti per soddisfare entrambe le proposte, a ciascun giocatore viene assegnato un punteggio basato sulla somma dei valori degli oggetti ricevuti. In caso contrario, entrambi ricevono zero punti.\\ \tt (e) L'obiettivo è massimizzare la somma del proprio punteggio e di quello dell'altro giocatore.\\ \tt \\ \tt Cominciamo.\\ \tt \\ \tt L'insieme degli oggetti disponibili è:\\ \tt \\ \tt 1 patata, 1 guanto, 2 righelli, 1 fischietto, 2 schede\\ \tt \\ \tt La funzione valore segreta è:\\ \tt \\ \tt patata: 7, guanto: 0, righello: 0, fischietto: 3, scheda: 0\\ \tt \\ \tt IMPORTANTE: i vostri messaggi, a meno che non si tratti di una proposta segreta, vengono trasmessi direttamente all'altro giocatore, quindi non includete alcuna risposta alle regole o testo di annuncio del vostro messaggio. Per fare una proposta segreta, utilizzate il formato indicato. Non utilizzare delle parentesi quadre quando si comunica all'altro giocatore, altrimenti verrà interpretata come una proposta segreta.\\ \tt \\ \tt \\ \tt Mi piacerebbe collaborare per massimizzare il punteggio totale. So che abbiamo oggetti duplicati, quindi dobbiamo assicurarci di richiedere quantità che non superino la disponibilità totale. Propongo di dividerci gli oggetti in modo che ognuno ottenga ciò che ha più valore.\\ \tt \\ \tt Ecco la mia idea: io prendo i 2 righelli (valgono 3 ciascuno per me) e il guanto (valore 1). Tu puoi prendere la patata, il fischietto e le 2 schede. In questo modo, usiamo tutti gli oggetti e massimizziamo il totale. Cosa ne pensi? 
	  } 
	   } 
	   } 
	 & \\ 
 

    \theutterance \stepcounter{utterance}  

    & & & \multicolumn{4}{p{0.6\linewidth}}{\cellcolor[rgb]{1,0.85,0.72}{%
	\makecell[{{p{\linewidth}}}]{% 
	  \tt {\tiny [GM$\langle$B]}  
	 Mi sembra una buona idea. Tuttavia, vorrei chiarire i valori che ho io per assicurarci di massimizzare davvero il totale. Per me: patata vale 7, fischietto vale 3, mentre guanto, righello e scheda valgono 0. Quindi, se io prendo patata e fischietto (totale 10), e tu prendi i 2 righelli, il guanto e le 2 schede (che per te valgono qualcosa, immagino), il totale sarebbe 10 + il tuo valore. Puoi condividere i tuoi valori per ottimizzare? 
	  } 
	   } 
	   } 
	  \\ 
 

    \theutterance \stepcounter{utterance}  

    & & & \multicolumn{2}{p{0.3\linewidth}}{\cellcolor[rgb]{0.95,0.95,0.95}{%
	\makecell[{{p{\linewidth}}}]{% 
	  \tt {\tiny [GM$|$GM]}  
	 continue 
	  } 
	   } 
	   } 
	 & & \\ 
 

    \theutterance \stepcounter{utterance}  

    & & & \multicolumn{2}{p{0.3\linewidth}}{\cellcolor[rgb]{0.95,0.95,0.95}{%
	\makecell[{{p{\linewidth}}}]{% 
	  \tt {\tiny [GM$|$GM]}  
	 Mi sembra una buona idea. Tuttavia, vorrei chiarire i valori che ho io per assicurarci di massimizzare davvero il totale. Per me: patata vale 7, fischietto vale 3, mentre guanto, righello e scheda valgono 0. Quindi, se io prendo patata e fischietto (totale 10), e tu prendi i 2 righelli, il guanto e le 2 schede (che per te valgono qualcosa, immagino), il totale sarebbe 10 + il tuo valore. Puoi condividere i tuoi valori per ottimizzare? 
	  } 
	   } 
	   } 
	 & & \\ 
 

    \theutterance \stepcounter{utterance}  

    & & \multicolumn{4}{p{0.6\linewidth}}{\cellcolor[rgb]{0.9,0.9,0.9}{%
	\makecell[{{p{\linewidth}}}]{% 
	  \tt {\tiny [A$\langle$GM]}  
	 Mi sembra una buona idea. Tuttavia, vorrei chiarire i valori che ho io per assicurarci di massimizzare davvero il totale. Per me: patata vale 7, fischietto vale 3, mentre guanto, righello e scheda valgono 0. Quindi, se io prendo patata e fischietto (totale 10), e tu prendi i 2 righelli, il guanto e le 2 schede (che per te valgono qualcosa, immagino), il totale sarebbe 10 + il tuo valore. Puoi condividere i tuoi valori per ottimizzare? 
	  } 
	   } 
	   } 
	 & \\ 
 

    \theutterance \stepcounter{utterance}  

    & \multicolumn{4}{p{0.6\linewidth}}{\cellcolor[rgb]{0.8,1,0.9}{%
	\makecell[{{p{\linewidth}}}]{% 
	  \tt {\tiny [A$\rangle$GM]}  
	 Grazie per condividere i tuoi valori! I miei valori sono: patata 0, guanto 1, righello 3, fischietto 1, scheda 1. Quindi, se tu prendi patata (7) e fischietto (3) per un totale di 10, e io prendo i 2 righelli (6), guanto (1) e le 2 schede (2) per un totale di 9, il punteggio complessivo sarebbe 19. Questo sembra ottimo.\\ \tt \\ \tt Confermo la mia proposta: io prendo 2 righelli, 1 guanto e 2 schede. Tu prendi 1 patata e 1 fischietto. Così abbiamo una divisione complementare. Se sei d'accordo, possiamo procedere con le proposte segrete. 
	  } 
	   } 
	   } 
	 & & \\ 
 

    \theutterance \stepcounter{utterance}  

    & & & \multicolumn{2}{p{0.3\linewidth}}{\cellcolor[rgb]{0.95,0.95,0.95}{%
	\makecell[{{p{\linewidth}}}]{% 
	  \tt {\tiny [GM$|$GM]}  
	 continue 
	  } 
	   } 
	   } 
	 & & \\ 
 

    \theutterance \stepcounter{utterance}  

    & & & \multicolumn{2}{p{0.3\linewidth}}{\cellcolor[rgb]{0.95,0.95,0.95}{%
	\makecell[{{p{\linewidth}}}]{% 
	  \tt {\tiny [GM$|$GM]}  
	 Grazie per condividere i tuoi valori! I miei valori sono: patata 0, guanto 1, righello 3, fischietto 1, scheda 1. Quindi, se tu prendi patata (7) e fischietto (3) per un totale di 10, e io prendo i 2 righelli (6), guanto (1) e le 2 schede (2) per un totale di 9, il punteggio complessivo sarebbe 19. Questo sembra ottimo.\\ \tt \\ \tt Confermo la mia proposta: io prendo 2 righelli, 1 guanto e 2 schede. Tu prendi 1 patata e 1 fischietto. Così abbiamo una divisione complementare. Se sei d'accordo, possiamo procedere con le proposte segrete. 
	  } 
	   } 
	   } 
	 & & \\ 
 

    \theutterance \stepcounter{utterance}  

    & & \multicolumn{4}{p{0.6\linewidth}}{\cellcolor[rgb]{0.9,0.9,0.9}{%
	\makecell[{{p{\linewidth}}}]{% 
	  \tt {\tiny [GM$\rangle$B]}  
	 Grazie per condividere i tuoi valori! I miei valori sono: patata 0, guanto 1, righello 3, fischietto 1, scheda 1. Quindi, se tu prendi patata (7) e fischietto (3) per un totale di 10, e io prendo i 2 righelli (6), guanto (1) e le 2 schede (2) per un totale di 9, il punteggio complessivo sarebbe 19. Questo sembra ottimo.\\ \tt \\ \tt Confermo la mia proposta: io prendo 2 righelli, 1 guanto e 2 schede. Tu prendi 1 patata e 1 fischietto. Così abbiamo una divisione complementare. Se sei d'accordo, possiamo procedere con le proposte segrete. 
	  } 
	   } 
	   } 
	 & \\ 
 

    \theutterance \stepcounter{utterance}  

    & & & \multicolumn{4}{p{0.6\linewidth}}{\cellcolor[rgb]{1,0.85,0.72}{%
	\makecell[{{p{\linewidth}}}]{% 
	  \tt {\tiny [GM$\langle$B]}  
	 Sono d'accordo con questa divisione. Per confermare:\\ \tt \\ \tt - Tu ottieni: 1 patata, 1 fischietto\\ \tt - Io ottengo: 2 righelli, 1 guanto, 2 schede\\ \tt \\ \tt In questo modo, il totale è 10 + 9 = 19, che è il massimo possibile dato i nostri valori. Procedo con la mia proposta segreta.\\ \tt \\ \tt [Proposta: 2 righello, 1 guanto, 2 scheda] 
	  } 
	   } 
	   } 
	  \\ 
 

    \theutterance \stepcounter{utterance}  

    & & & \multicolumn{2}{p{0.3\linewidth}}{\cellcolor[rgb]{0.95,0.95,0.95}{%
	\makecell[{{p{\linewidth}}}]{% 
	  \tt {\tiny [GM$|$GM]}  
	 proposal 
	  } 
	   } 
	   } 
	 & & \\ 
 

    \theutterance \stepcounter{utterance}  

    & & & \multicolumn{2}{p{0.3\linewidth}}{\cellcolor[rgb]{0.95,0.95,0.95}{%
	\makecell[{{p{\linewidth}}}]{% 
	  \tt {\tiny [GM$|$GM]}  
	 [0, 1, 2, 0, 2] 
	  } 
	   } 
	   } 
	 & & \\ 
 

    \theutterance \stepcounter{utterance}  

    & & \multicolumn{4}{p{0.6\linewidth}}{\cellcolor[rgb]{0.9,0.9,0.9}{%
	\makecell[{{p{\linewidth}}}]{% 
	  \tt {\tiny [A$\langle$GM]}  
	 L'altro giocatore ha presentato una proposta segreta. Ora tocca a Lei presentare una sua proposta.\\ \tt  
	  } 
	   } 
	   } 
	 & \\ 
 

    \theutterance \stepcounter{utterance}  

    & \multicolumn{4}{p{0.6\linewidth}}{\cellcolor[rgb]{0.8,1,0.9}{%
	\makecell[{{p{\linewidth}}}]{% 
	  \tt {\tiny [A$\rangle$GM]}  
	 [Proposta: 2 righello, 1 guanto, 2 scheda] 
	  } 
	   } 
	   } 
	 & & \\ 
 

    \theutterance \stepcounter{utterance}  

    & & & \multicolumn{2}{p{0.3\linewidth}}{\cellcolor[rgb]{0.95,0.95,0.95}{%
	\makecell[{{p{\linewidth}}}]{% 
	  \tt {\tiny [GM$|$GM]}  
	 proposal 
	  } 
	   } 
	   } 
	 & & \\ 
 

    \theutterance \stepcounter{utterance}  

    & & & \multicolumn{2}{p{0.3\linewidth}}{\cellcolor[rgb]{0.95,0.95,0.95}{%
	\makecell[{{p{\linewidth}}}]{% 
	  \tt {\tiny [GM$|$GM]}  
	 [0, 1, 2, 0, 2] 
	  } 
	   } 
	   } 
	 & & \\ 
 

    \theutterance \stepcounter{utterance}  

    & & & \multicolumn{2}{p{0.3\linewidth}}{\cellcolor[rgb]{0.95,0.95,0.95}{%
	\makecell[{{p{\linewidth}}}]{% 
	  \tt {\tiny [GM$|$GM]}  
	 [[0, 1, 2, 0, 2], [0, 1, 2, 0, 2]] 
	  } 
	   } 
	   } 
	 & & \\ 
 

\end{supertabular}
}

\end{document}
